\newif\ifglobal

%%%%%%%%%%%%%%%%%%%%%%%
%%% Compile to view %%%
%%%%%%%%%%%%%%%%%%%%%%%

%%%%%%%%%%%%%%%%%%%%%%%%%%%%%%%%%%%%%%%%%%%%%%%%%%%
%%% Readme:                                     %%%
%%% https://www.overleaf.com/read/bwdjvfjksvjy/ %%%
%%%%%%%%%%%%%%%%%%%%%%%%%%%%%%%%%%%%%%%%%%%%%%%%%%%

% >>> M?: marks a module setting number or important notice

% >>> M4: Set {./relative-path} to external files for this module <<<
\def \modulefiles {./17-TIMG}
% To add an external file, use:
% (Replace '---' with actual filename)
% '\input{\modulefiles/tables/---.tex}' for tables
% '\includegraphics{\modulefiles/figures/---}' for figures
% '\subfile{\modulefiles/---}' for register files


% >>> M1: This module is in global project? <<<
% 'YES' - uncomment; 'NO' - comment out
\globaltrue

\ifglobal
    \documentclass[main\_\_CN.tex]{subfiles}
   
\else
    % Font "FandolSong-Regular" used by ctex does not have GSUB table.
    % It causes warnings that do not affect compilation and can be ignored.
    \PassOptionsToPackage{quiet}{fontspec}

    \documentclass[openany, 10pt]{book}
   
    % >>> M2: In ./00-shared/preamble-trm-module, also find
    %         settings for visibility of todo notes, watermark <<< 
    \usepackage{preamble-trm-module}

    % Enable Chinese typesetting
    \usepackage{ctex}

    % >>> M3: Temporary labels in Chinese
    %         you can add your own labels there <<<
    \myexternaldocument{./00-shared/temp-labels-cn}

\fi %\ifglobal


\setcounter{secnumdepth}{4}


% >>> M5: If this module needs any updates in
% './00-shared/preamble-trm-module.sty'
% (include additional package, etc.), contact your module PM
% or create the file './00-shared/changelog.tex' make the
% necessary updates yourself and document them in this file <<<

% Make text left-aligned and allow latex to break pages naturally, comment out for full justification
\raggedright
\raggedbottom


\begin{document}


% Sets language into which pre-defined bits of text are translated
\selectlanguage{chinese}
\setchineseformatting


% Do not delete this block
\ifglobal
% Add the following front matter if
% this module is outside of global project
\else
    \InputIfFileExists{readme}{}{}
    \clearpage
    \listoftodos
    \listofdonetodos
    \tableofcontents
    %\listoftables
    %\listoffigures
\fi


%%%%%%%%%%%%%%%%%%%%%%%%%
%%% TRM Module Begins %%%
%%%%%%%%%%%%%%%%%%%%%%%%%

\hypertarget{timg}{}
\chapter{定时器组 (TIMG)}
\label{mod:timg}

%%%%%%%%%%%%%%%%
%%% Overview %%%
%%%%%%%%%%%%%%%%

\section{概述}
\label{sec:timg-overview}

通用定时器可用于准确设定时间间隔、在一定间隔后触发（周期或非周期的）中断或充当硬件时钟。如图 \ref{fig:timg-overview} \textit{\nameref{fig:timg-overview}} 所示，\chipname{} 包含两个定时器组，即定时器组 0 和定时器组 1（下文用 TIMG\regindex{n} 表示，\regindex{n} 为 0 或 1）。每个定时器组有两个通用定时器（下文用 T\regindex{x} 表示，\regindex{x} 为 0 或 1）和一个主系统看门狗定时器。此外，TIMG0 还包含一个使用 XTAL\_CLK 作为基准时钟的频率计算模块 FC (Frequency Calculation)，详细参考第 \ref{sec:rtc-calibration-desc} 节 \textit{\nameref{sec:rtc-calibration-desc}}。通用定时器基于 16 位预分频器和 54 位可自动重新加载的可逆计数器。

\begin{figure}[H]
    \centering
    \includegraphics[width=0.46\textwidth]{\modulefiles/figures/timg_units}
    \caption{定时器组概览}
    \label{fig:timg-overview}
\end{figure}

本章虽然包含主系统看门狗定时器的寄存器描述，但其功能描述请参阅章节 \ref{mod:wdt} \textit{\nameref{mod:wdt}}。本章中“定时器”指代通用定时器。

%%%%%%%%%%%%%%%%%%%%
%%% Feature List %%%
%%%%%%%%%%%%%%%%%%%%

\section{主要特性}
\label{sec:timg-features}

定时器具有如下功能：

\begin{itemize}
    \item 三个时钟源：REF\_80M\_CLK、XTAL\_CLK、FOSC\_20M\_CLK 时钟
    \item 16 位时钟预分频器，分频系数为 2 到 65536
    \item 54 位时基计数器，可配置成递增或递减
    \item 可读取时基计数器的实时值
    \item 暂停和恢复时基计数器
    \item 可配置的警报生成机制
    \item 计数器值重新加载 —— 报警时自动重新加载或软件控制的即时重新加载
    \item TIMG0 定时器组慢速时钟频率计算
    \item 电平触发中断
    \item 支持多个 ETM 任务和事件
\end{itemize}

%%%%%%%%%%%%%%%%%%%%%%%%%%%%%%
%%% Architectural Overview %%%
%%%%%%%%%%%%%%%%%%%%%%%%%%%%%%

\section{结构概览}
\label{sec:timg-arch-overview}

\begin{figure}[H]
    \centering
    \includegraphics[width=0.9\textwidth]{\modulefiles/figures/timg-architecture}
    \caption{定时器组架构}
    \label{fig:timg-architecture}
\end{figure}

图 \ref{fig:timg-architecture} \textit{\nameref{fig:timg-architecture}} 为定时器组的 T\regindex{x}。T\regindex{x} 包含一个 16 位整数预分频器、一个时基计数器和一个用于产生警报的比较器。

%%%%%%%%%%%%%%%%%%%%%%%%%%%%%%
%%% Functional Description %%%
%%%%%%%%%%%%%%%%%%%%%%%%%%%%%%

\section{功能描述}
\label{sec:timg-func-descr}

\subsection{16 位预分频器与时钟选择器}

\begin{itemize}
    \item 定时器组 TIMG\regindex{n} 的定时器 T\regindex{x} 可通过配置 T\regindex{x}\_SRC\_SEL 来选择时钟源，T\regindex{x}\_SRC\_SEL 是寄存器 \hyperref[regdesc:HPSYSCLKRSTTIMERGRPNCTRL0REG]{HP\_SYS\_CLKRST\_TIMERGRP\regindex{n}\_CTRL0\_REG} 的 \hyperref[fielddesc:HPSYSCLKRSTTIMERGRPNTXSRCSEL]{HP\_SYS\_CLKRST\_TIMERGRP\regindex{n}\_T\regindex{x}\_SRC\_SEL} 字段。
    \begin{itemize}
        \item 该字段为 0 时，选择 XTAL\_CLK。
        \item 为 1 时，选择 FOSC\_20M\_CLK。
        \item 为 2 时，选择 REF\_80M\_CLK。
    \end{itemize}

    \item 所选时钟源需通过 T\regindex{x}\_CLK\_EN 使能，T\regindex{x}\_CLK\_EN 是 \hyperref[regdesc:HPSYSCLKRSTTIMERGRPNCTRL0REG]{HP\_SYS\_CLKRST\_TIMERGRP\regindex{n}\_CTRL0\_REG} 寄存器的 \hyperref[fielddesc:HPSYSCLKRSTTIMERGRPNTXCLKEN]{HP\_SYS\_CLKRST\_TIMERGRP\regindex{n}\_T\regindex{x}\_CLK\_EN} 字段（见图 \ref{fig:timg-architecture} \textit{\nameref{fig:timg-architecture}}）。置位该字段可开启所选时钟，清零则关闭。
    \item 启用的时钟信号将输入至 16 位预分频器，经分频后产生时基计数器所使用的时基计数器时钟 (TB\_CLK)。16 位预分频器的分频系数可通过 \hyperref[fielddesc:TIMGTXDIVIDER]{TIMG\_T\regindex{x}\_DIVIDER} 字段配置。\hyperref[fielddesc:TIMGTXDIVIDER]{TIMG\_T\regindex{x}\_DIVIDER} 字段可配置为 0 \verb+~+ 65535 之间的任意值，分频系数的范围为 2 \verb+~+ 65536。具体如下：
    \begin{itemize}
        \item 字段置 0 时：分频系数为 65536
        \item 字段置 1 时：分频系数为 2
        \item 字段置 2 时：分频系数仍为 2
        \item 字段置 3 \verb+~+ 65535 时：分频系数为 3 \verb+~+ 65535
    \end{itemize}

    若需更改 16 位预分频器，需先重新配置 \hyperref[fielddesc:TIMGTXDIVIDER]{TIMG\_T\regindex{x}\_DIVIDER} 字段，再将 \hyperref[fielddesc:TIMGTXDIVCNTRST]{TIMG\_T\regindex{x}\_DIVCNT\_RST} 置 1，同时需要关闭定时器（即 \hyperref[fielddesc:TIMGTXEN]{TIMG\_T\regindex{x}\_EN} 必须清零），否则会造成不可预知的结果。

\end{itemize}

\subsection{54 位时基计数器}

54 位时基计数器基于 TB\_CLK，具有以下功能：

\begin{itemize}
    \item 计数方向可通过 \hyperref[fielddesc:TIMGTXINCREASE]{TIMG\_T\regindex{x}\_INCREASE} 字段配置，可设置为递增或递减。
    \item 可通过置位或清零 \hyperref[fielddesc:TIMGTXEN]{TIMG\_T\regindex{x}\_EN} 字段使能或关闭计数器：
    \begin{itemize}
        \item 使能时，计数器的值会在每个 TB\_CLK 周期递增或递减。
        \item 关闭时，计数器暂停计数。
        \item 注意，无论 \hyperref[fielddesc:TIMGTXEN]{TIMG\_T\regindex{x}\_EN} 是否置位，\hyperref[fielddesc:TIMGTXINCREASE]{TIMG\_T\regindex{x}\_INCREASE} 字段均可修改，计数方向会立即生效。
    \end{itemize}

    \item 由于 CPU 为 32 位架构，读取 54 位计数值需通过锁存机制完成：
    \begin{itemize}
    \item 向 \hyperref[regdesc:TIMGTXUPDATEREG]{TIMG\_T\regindex{x}UPDATE\_REG} 写入任意值以触发锁存。
    \item 当前计数值会被锁存至寄存器 \hyperref[regdesc:TIMGTXLOREG]{TIMG\_T\regindex{x}LO\_REG}（低 32 位）和 \hyperref[regdesc:TIMGTXHIREG]{TIMG\_T\regindex{x}HI\_REG}（高 22 位）。
    \item 当 \hyperref[regdesc:TIMGTXUPDATEREG]{TIMG\_T\regindex{x}UPDATE\_REG} 被硬件清零，表示锁存完成，可以从上述寄存器中读取当前值。
    \item 在 \hyperref[regdesc:TIMGTXUPDATEREG]{TIMG\_T\regindex{x}UPDATE\_REG} 被写入新值之前，保持寄存器 \hyperref[regdesc:TIMGTXLOREG]{TIMG\_T\regindex{x}LO\_REG} 和 \hyperref[regdesc:TIMGTXHIREG]{TIMG\_T\regindex{x}HI\_REG} 的值不变，确保 CPU 可稳定读取。
    \end{itemize}
\end{itemize}


\subsection{警报生成}

定时器可配置为在当前值与报警值相同时触发报警。报警会产生中断，同时（可选择）让定时器的当前值自动重新加载（详见第 \ref{sec:timg-func-descr-reload} 节 \textit{\nameref{sec:timg-func-descr-reload}}）。

54 位报警值可在 \hyperref[regdesc:TIMGTXALARMLOREG]{TIMG\_T\regindex{x}ALARMLO\_REG} 和 \hyperref[regdesc:TIMGTXALARMHIREG]{TIMG\_T\regindex{x}ALARMHI\_REG} 配置，两者分别代表报警值的低 32 位和高 22 位。但是，只有置位 \hyperref[fielddesc:TIMGTXALARMEN]{TIMG\_T\regindex{x}\_ALARM\_EN} 字段使能报警功能后，配置的报警值才会生效。配置报警需要多次写寄存器，在此过程中计数器持续计数，配置完成时定时器当前值可能已经越过目标报警值。为避免这种情况下报警始终不触发，置位 \hyperref[fielddesc:TIMGTXALARMEN]{TIMG\_T\regindex{x}\_ALARM\_EN} 时，若定时器当前值出现以下情况，硬件会立即触发报警：

\begin{itemize}
    \item 可逆计数器向上计数时，定时器的当前值高于报警值（在一定范围内）
    \item 可逆计数器向下计数时，定时器的当前值低于报警值（在一定范围内）
\end{itemize}

表 \ref{tab:timg-alarm-generation-up} \textit{\nameref{tab:timg-alarm-generation-up}} 和表 \ref{tab:timg-alarm-generation-down} \textit{\nameref{tab:timg-alarm-generation-down}} 说明了定时器当前值、报警值与报警触发的关系。假设定时器当前值和报警值如下：

\begin{itemize}
    \item TIMG\_VALUE = \{\hyperref[regdesc:TIMGTXHIREG]{TIMG\_T\regindex{x}HI\_REG}, \hyperref[regdesc:TIMGTXLOREG]{TIMG\_T\regindex{x}LO\_REG}\}
    \item ALARM\_VALUE = \{\hyperref[regdesc:TIMGTXALARMHIREG]{TIMG\_T\regindex{x}ALARMHI\_REG}, \hyperref[regdesc:TIMGTXALARMLOREG]{TIMG\_T\regindex{x}ALARMLO\_REG}\}
\end{itemize}

\begin{table}[H]
\centering
\caption{可逆计数器向上计数时的报警场景}
\label{tab:timg-alarm-generation-up}

\begin{tabular}{|l|l|m{7cm}|}
\hline
\rowcolor{lightgray}
    \textbf{场景} & \textbf{范围} & \textbf{报警} \\\hline
     1 & ALARM\_VALUE $-$ TIMG\_VALUE $>$ $2^{53}$ & 触发 \\\hline
     2 & 0 $<$ ALARM\_VALUE $-$ TIMG\_VALUE $\leq$ $2^{53}$ & 可逆计数器向上计数，TIMG\_VALUE 达到 ALARM\_VALUE 时报警 \\\hline
     3 & 0 $\leq$ TIMG\_VALUE $-$ ALARM\_VALUE $<$ $2^{53}$ & 触发 \\\hline
     4 & TIMG\_VALUE $-$ ALARM\_VALUE $\geq$ $2^{53}$ & 可逆计数器向上计数达到最大值时，重新开始从 0 向上计数，TIMG\_VALUE 达到 ALARM\_VALUE 时触发报警\\\hline
\end{tabular}
\end{table}

\begin{table}[H]
\centering
\caption{可逆计数器向下计数时的报警场景}
\label{tab:timg-alarm-generation-down}

\begin{tabular}{|l|l|m{7cm}|}
\hline
\rowcolor{lightgray}
    \textbf{场景} & \textbf{范围} & \textbf{报警} \\\hline
     5 & TIMG\_VALUE $-$ ALARM\_VALUE  $>$ $2^{53}$ & 触发 \\\hline
     6 &  0 $<$ TIMG\_VALUE $-$ ALARM\_VALUE $\leq$ $2^{53}$ & 可逆计数器向下计数，TIMG\_VALUE 达到 ALARM\_VALUE 时报警 \\\hline
     7 & 0 $\leq$ ALARM\_VALUE $-$ TIMG\_VALUE $<$ $2^{53}$ & 触发 \\\hline
     8 & ALARM\_VALUE $-$ TIMG\_VALUE $\geq$ $2^{53}$ & 可逆计数器向下计数达到最小值时，重新开始从最大值向下计数，TIMG\_VALUE 达到 ALARM\_VALUE 时触发报警 \\\hline
\end{tabular}
\end{table}

报警时，\hyperref[fielddesc:TIMGTXALARMEN]{TIMG\_T\regindex{x}\_ALARM\_EN} 字段自动清零，在下次置位 \hyperref[fielddesc:TIMGTXALARMEN]{TIMG\_T\regindex{x}\_ALARM\_EN} 前不会再次报警。

\subsection{定时器重新加载}
\label{sec:timg-func-descr-reload}

定时器重新加载指将定时器的低 32 位和高 22 位分别更新为 \hyperref[fielddesc:TIMGTXLOADLO]{TIMG\_T\regindex{x}\_LOAD\_LO} 和 \hyperref[fielddesc:TIMGTXLOADHI]{TIMG\_T\regindex{x}\_LOAD\_HI} 字段中存储的值。但仅写入这两个字段不会立即改变定时器当前值，写入的重新加载值会被定时器忽视，直到重新加载事件触发后才生效。重新加载事件可通过以下两种方式触发：

\begin{itemize}
    \item \textbf{软件触发即时重新加载：}CPU 向寄存器 \hyperref[regdesc:TIMGTXLOADREG]{TIMG\_T\regindex{x}LOAD\_REG} 写入任意值时，立即触发重新加载事件。
    \begin{itemize}
        \item 若 \hyperref[fielddesc:TIMGTXEN]{TIMG\_T\regindex{x}\_EN} 为置位状态，定时器会从新数值开始继续递增或递减计数。若此时 \hyperref[fielddesc:TIMGTXALARMEN]{TIMG\_T\regindex{x}\_ALARM\_EN} 置位，仍会根据表 \ref{tab:timg-alarm-generation-up} \textit{\nameref{tab:timg-alarm-generation-up}} 或表 \ref{tab:timg-alarm-generation-down} \textit{\nameref{tab:timg-alarm-generation-down}} 中的规则在对应时刻产生报警。
        \item 若 \hyperref[fielddesc:TIMGTXEN]{TIMG\_T\regindex{x}\_EN} 是清零状态，定时器将保持当前值，直至计数重新使能。
    \end{itemize}
    \item \textbf{报警触发自动重新加载：}若启用了自动重新加载功能，即 \hyperref[fielddesc:TIMGTXAUTORELOAD]{TIMG\_T\regindex{x}\_AUTORELOAD} 字段为 1，定时器将在报警时自动重新加载，从重新加载值开始继续递增或递减计数。该功能常用于周期性报警时重置定时器的值。若未启用该功能，报警后定时器的值会在过报警值后继续递增或递减，不会自动重新加载。
\end{itemize}


\subsection{TIMG0 定时器组慢速时钟频率计算}
\label{sec:rtc-calibration-desc}

TIMG0 定时器组可以以 XTAL\_CLK 为基准时钟，计算 \chipname{} 提供的慢速时钟的实际频率，输入给定时器组的具体慢速时钟解释请查看第 \ref{sec:rtc-calibration-procedure} 节 \textit{\nameref{sec:rtc-calibration-procedure}}。计算方式如下：

\begin{enumerate}
    \item 通过周期性或单次计算的方式启动频率计算模块（详见第 \ref{sec:rtc-calibration-procedure} 节 \textit{\nameref{sec:rtc-calibration-procedure}}）。
    \item 在接收到计算开始的信号后，两个分别工作在 XTAL\_CLK 以及慢速时钟的计数器同时开始计数，当慢速时钟的计数器达到设定的计算周期 C0 时，同时停止两个计数器。
    \item 通过 XTAL\_CLK 的计数器值 C1 即可计算慢速时钟的时钟频率：$f\_rtc = \frac{C0 \times f\_XTAL\_CLK}{C1}$。
\end{enumerate}

请注意慢速时钟频率计算功能使用时应保证输入频率不大于 $\frac{f\_XTAL\_CLK}{2}$。


\subsection{复位}


TIMG\regindex{n} 支持完全复位，不支持部分复位，即没有单独复位 TIMG\regindex{n} 模块中某个组件的复位。复位可通过向相应字段先写 1 后写 0 实现：

\begin{itemize}
    \item 完全复位：
        \begin{itemize}
            \item TIMG0：\hyperref[fielddesc:HPSYSCLKRSTTIMERGRP0RSTEN]{HP\_SYS\_CLKRST\_TIMERGRP0\_RST\_EN}
            \item TIMG1：\hyperref[fielddesc:HPSYSCLKRSTTIMERGRP1RSTEN]{HP\_SYS\_CLKRST\_TIMERGRP1\_RST\_EN}
        \end{itemize}
    \item 部分复位：不支持
\end{itemize}

更多信息可参考章节 \ref{mod:resclk} \textit{\nameref{mod:resclk}}。


%%%%%%%%%%%%%%%%%%%%%%%%%%%%%%%%%
%%% Event Task Matrix Feature %%%
%%%%%%%%%%%%%%%%%%%%%%%%%%%%%%%%%

\section{事件任务矩阵功能}
\label{sec:timg-etm-func}

在 \chipname{} 中，定时器组支持事件任务矩阵 (ETM) 功能，即可以通过任意外设的 ETM 事件触发定时器组的 ETM 任务，或者通过定时器组的 ETM 事件触发任意外设的 ETM 任务。关于 ETM 更多详细信息，请参考章节 \ref{mod:evntaskmatrix} \textit{\nameref{mod:evntaskmatrix}}。这里仅介绍与定时器组相关的 ETM 任务和 ETM 事件。

定时器组可接收的 ETM 任务有：

\begin{itemize}
\item TG0\_TASK\_CNT\_START\_TIMER\regindex{x}：触发时启用定时器组 0 中 T\regindex{x} 的时基计数器
\item TG1\_TASK\_CNT\_START\_TIMER\regindex{x}：触发时启用定时器组 1 中 T\regindex{x} 的时基计数器
\item TG0\_TASK\_CNT\_STOP\_TIMER\regindex{x}：触发时关闭定时器组 0 中 T\regindex{x} 的时基计数器
\item TG1\_TASK\_CNT\_STOP\_TIMER\regindex{x}：触发时关闭定时器组 1 中 T\regindex{x} 的时基计数器

\begin{tiplisting}
以上两个 ETM 任务的功能，与通过 APB（高级外设总线）配置 \hyperref[fielddesc:TIMGTXEN]{TIMG\_T\regindex{x}\_EN} 字段实现的功能相同。当这些操作同时发生时，优先级由高到低如下：

\begin{enumerate}
    \item TG0\_TASK\_CNT\_START\_TIMER\regindex{x} 与 TG1\_TASK\_CNT\_START\_TIMER\regindex{x}：触发时启用时基计数器
    \item TG0\_TASK\_CNT\_STOP\_TIMER\regindex{x} 与 TG1\_TASK\_CNT\_STOP\_TIMER\regindex{x}：触发时关闭时基计数器
    \item APB 方式配置 \hyperref[fielddesc:TIMGTXEN]{TIMG\_T\regindex{x}\_EN}：启用或关闭时基计数器
\end{enumerate}
\end{tiplisting}

\item TG0\_TASK\_ALARM\_START\_TIMER\regindex{x}：触发时启用定时器组 0 中 T\regindex{x} 的警报生成
\item TG1\_TASK\_ALARM\_START\_TIMER\regindex{x}：触发时启用定时器组 1 中 T\regindex{x} 的警报生成

\begin{tiplisting}
除了 ETM 任务事件触发外，通过报警中断或通过 APB 方式配置 \hyperref[fielddesc:TIMGTXALARMEN]{TIMG\_T\regindex{x}\_ALARM\_EN} 也可以控制警报生成。当这些操作同时发生时，其优先级从高到低依次为：


\begin{enumerate}
    \item TG0\_TASK\_ALARM\_START\_TIMER\regindex{x} 与 TG1\_TASK\_ALARM\_START\_TIMER\regindex{x}：触发时启用警报生成
    \item 报警中断：触发时关闭警报生成
    \item APB 方式配置 \hyperref[fielddesc:TIMGTXALARMEN]{TIMG\_T\regindex{x}\_ALARM\_EN}：启用/关闭警报生成
\end{enumerate}
\end{tiplisting}

\item TG0\_TASK\_CNT\_CAP\_TIMER\regindex{x}：触发时将定时器组 0 中 T\regindex{x} 的当前计数值更新至 \hyperref[regdesc:TIMGTXLOREG]{TIMG\_T\regindex{x}LO\_REG} 和 \hyperref[regdesc:TIMGTXHIREG]{TIMG\_T\regindex{x}HI\_REG} 寄存器
\item TG1\_TASK\_CNT\_CAP\_TIMER\regindex{x}：触发时将定时器组 1 中 T\regindex{x} 的当前计数值更新至 \hyperref[regdesc:TIMGTXLOREG]{TIMG\_T\regindex{x}LO\_REG} 和 \hyperref[regdesc:TIMGTXHIREG]{TIMG\_T\regindex{x}HI\_REG} 寄存器
\item TG0\_TASK\_CNT\_RELOAD\_TIMER\regindex{x}：触发时用存储在 \hyperref[fielddesc:TIMGTXLOADLO]{TIMG\_T\regindex{x}\_LOAD\_LO} 和 \hyperref[fielddesc:TIMGTXLOADHI]{TIMG\_T\regindex{x}\_LOAD\_HI} 中的重新加载值覆盖定时器组 0 中 T\regindex{x} 的当前计数值
\item TG1\_TASK\_CNT\_RELOAD\_TIMER\regindex{x}：触发时用存储在 \hyperref[fielddesc:TIMGTXLOADLO]{TIMG\_T\regindex{x}\_LOAD\_LO} 和 \hyperref[fielddesc:TIMGTXLOADHI]{TIMG\_T\regindex{x}\_LOAD\_HI} 中的重新加载值覆盖定时器组 1 中 T\regindex{x} 的当前计数值
\end{itemize}

定时器组可产生的 ETM 事件有：

\begin{itemize}
\item TG0\_EVT\_CNT\_CMP\_TIMER\regindex{x}：表示定时器组 0 中的 T\regindex{x} 定时器中断事件
\item TG1\_EVT\_CNT\_CMP\_TIMER\regindex{x}：表示定时器组 1 中的 T\regindex{x} 定时器中断事件
\end{itemize}

只有当 \hyperref[fielddesc:TIMGETMEN]{TIMG\_ETM\_EN} 设置为 1 时，ETM 任务和事件才会生效。

在具体应用中，定时器组的 ETM 事件可以用来触发定时器组的 ETM 任务。例如，TG0\_EVT\_CNT\_CMP\_TIMER\regindex{x} 与 TG1\_EVT\_CNT\_CMP\_TIMER\regindex{x} 事件分别可以触发 TG0\_TASK\_ALARM\_START\_TIMER\regindex{x} 与 TG1\_TASK\_ALARM\_START\_TIMER\regindex{x} 任务，从而实现周期性报警。具体配置步骤请参考第 \ref{sec:timg-conf-etm} 节 \textit{\nameref{sec:timg-conf-etm}}。

%%%%%%%%%%%%%%%%%%
%%% Interrupts %%%
%%%%%%%%%%%%%%%%%%

\section{中断}
\label{sec:timg-interrupt}

\chipname{} 的 TIMG\regindex{n} 可以生成以下中断信号并发送给\hyperref[mod:intmtrx]{中断矩阵}。

    \begin{itemize}
        \item TG\regindex{n}\_T0\_INTR
        \item TG\regindex{n}\_T1\_INTR
        \item TG\regindex{n}\_WDT\_INTR
    \end{itemize}

这些中断信号由 TIMG\regindex{n} 的内部中断源生成。表 \ref{tab:timg-int-sources} \textit{\nameref{tab:timg-int-sources}} 列出了 TIMG\regindex{n} 的内部中断源、中断触发条件、以及生成的中断信号。

\begin{longtable}{ | L{5cm} | L{6.5cm} | L{4cm} | }
\caption{TIMG\regindex{n} 的内部中断源}
\label{tab:timg-int-sources}
\\ \hline
\rowcolor{lightgray}
\textbf{内部中断源} & \textbf{触发条件}  &  \textbf{中断信号}\\ \hline
\label{int:TIMGT0INT}\hyperref[int:TIMGT0INT]{TIMG\_T0\_INT}  &  T0 定时器产生警报 &  TG\regindex{n}\_T0\_INTR  \\ \hline
\label{int:TIMGT1INT}\hyperref[int:TIMGT1INT]{TIMG\_T1\_INT}  &  T1 定时器产生警报 &  TG\regindex{n}\_T1\_INTR  \\ \hline
\label{int:TIMGWDTINT}\hyperref[int:TIMGWDTINT]{TIMG\_WDT\_INT} &  看门狗定时器产生警报 & TG\regindex{n}\_WDT\_INTR \\ \hline
\end{longtable}

\begin{tiplisting}
中断、中断源、中断信号等术语的解释以及它们之间的联系请参考章节 \ref{mod:intmtrx} \textit{\nameref{mod:intmtrx}} > \ref{sec:intmtx-int-terminology} \textit{\nameref{sec:intmtx-int-terminology}}。
\end{tiplisting}

每个中断源都有一组通用的配置寄存器，详见章节 \hyperref[interrupt-config-registers]{\textit{中断配置寄存器}}。具体寄存器名称请查看第 \ref{sec:timg-reg-summ} 节 \textit{\nameref{sec:timg-reg-summ}}。

%%%%%%%%%%%%%%%%%%%%%%%%%%%%%%
%%% Programming Procedures %%%
%%%%%%%%%%%%%%%%%%%%%%%%%%%%%%

\section{配置流程}
\label{sec:timg-prog-procedures}

\subsection{定时器用作简单时钟}
\label{sec:timg-func-descr-simple-clock}
\begin{enumerate}
    \item 配置时基计数器
        \begin{itemize}
        \item 配置 \hyperref[fielddesc:HPSYSCLKRSTTIMERGRPNTXSRCSEL]{HP\_SYS\_CLKRST\_TIMERGRP\regindex{n}\_T\regindex{x}\_SRC\_SEL} 字段选择时钟源
        \item 配置 \hyperref[fielddesc:TIMGTXDIVIDER]{TIMG\_T\regindex{x}\_DIVIDER} 设置 16 位预分频器
        \item 置位或清除 \hyperref[fielddesc:TIMGTXINCREASE]{TIMG\_T\regindex{x}\_INCREASE} 配置定时器方向
        \item 在 \hyperref[fielddesc:TIMGTXLOADLO]{TIMG\_T\regindex{x}\_LOAD\_LO} 和 \hyperref[fielddesc:TIMGTXLOADHI]{TIMG\_T\regindex{x}\_LOAD\_HI} 上写初始值设置定时器的初始值，然后在 \hyperref[regdesc:TIMGTXLOADREG]{TIMG\_T\regindex{x}LOAD\_REG} 上写任意值将初始值重新加载进定时器
    \end{itemize}
    \item 置位 \hyperref[fielddesc:TIMGTXEN]{TIMG\_T\regindex{x}\_EN} 开启定时器
    \item 获得定时器的当前值
    \begin{itemize}
        \item 在 \hyperref[regdesc:TIMGTXUPDATEREG]{TIMG\_T\regindex{x}UPDATE\_REG} 上写任意值锁存定时器的当前值
        \item 等待硬件将 \hyperref[regdesc:TIMGTXUPDATEREG]{TIMG\_T\regindex{x}UPDATE\_REG} 清 0
        \item 从 \hyperref[regdesc:TIMGTXLOREG]{TIMG\_T\regindex{x}LO\_REG} 和 \hyperref[regdesc:TIMGTXHIREG]{TIMG\_T\regindex{x}HI\_REG} 读取锁存的定时器值
    \end{itemize}
\end{enumerate}

\subsection{定时器用于单次报警}
\label{sec:timg-func-descr-timer-alarm}

\begin{enumerate}
    \item 按照第 \ref{sec:timg-func-descr-simple-clock} 节 \textit{\nameref{sec:timg-func-descr-simple-clock}} 的第 1 步配置时基计数器
    \item 配置报警
    \begin{itemize}
        \item 设置 \hyperref[regdesc:TIMGTXALARMLOREG]{TIMG\_T\regindex{x}ALARMLO\_REG} 和 \hyperref[regdesc:TIMGTXALARMHIREG]{TIMG\_T\regindex{x}ALARMHI\_REG}，配置报警值
        \item 置位 \hyperref[fielddesc:TIMGTXINTENA]{TIMG\_T\regindex{x}\_INT\_ENA} 使能中断
    \end{itemize}
    \item 清零 \hyperref[fielddesc:TIMGTXAUTORELOAD]{TIMG\_T\regindex{x}\_AUTORELOAD} 关闭自动重新加载
    \item 置位 \hyperref[fielddesc:TIMGTXEN]{TIMG\_T\regindex{x}\_EN} 开启定时器
    \item 置位 \hyperref[fielddesc:TIMGTXALARMEN]{TIMG\_T\regindex{x}\_ALARM\_EN} 开启报警
    \item 处理报警中断
    \begin{itemize}
        \item 置位定时器在 \hyperref[fielddesc:TIMGTXINTCLR]{TIMG\_T\regindex{x}\_INT\_CLR} 的对应位清除中断
        \item 清零 \hyperref[fielddesc:TIMGTXEN]{TIMG\_T\regindex{x}\_EN} 关闭定时器
    \end{itemize}
\end{enumerate}

\subsection{通过 APB 设置定时器用于周期性报警}

\begin{enumerate}
    \item 按照第 \ref{sec:timg-func-descr-simple-clock} 节 \textit{\nameref{sec:timg-func-descr-simple-clock}} 的第 1 步配置时基计数器
    \item 按照第 \ref{sec:timg-func-descr-timer-alarm} 节 \textit{\nameref{sec:timg-func-descr-timer-alarm}} 的第 2 步配置报警
    \item 置位 \hyperref[fielddesc:TIMGTXAUTORELOAD]{TIMG\_T\regindex{x}\_AUTORELOAD} 使能自动重新加载，通过 \hyperref[fielddesc:TIMGTXLOADLO]{TIMG\_T\regindex{x}\_LOAD\_LO} 和 \hyperref[fielddesc:TIMGTXLOADHI]{TIMG\_T\regindex{x}\_LOAD\_HI} 配置重新加载值
    \item 置位 \hyperref[fielddesc:TIMGTXEN]{TIMG\_T\regindex{x}\_EN} 开启定时器
    \item 置位 \hyperref[fielddesc:TIMGTXALARMEN]{TIMG\_T\regindex{x}\_ALARM\_EN} 开启报警
    \item 处理报警中断（每次报警时重复）
    \begin{itemize}
        \item 置位定时器在 \hyperref[fielddesc:TIMGTXINTCLR]{TIMG\_T\regindex{x}\_INT\_CLR} 的对应位清除中断
        \item 如下一次报警需要新的报警值和重新加载值（即每次都有不同的报警间隔），则应根据需要重新配置 \hyperref[regdesc:TIMGTXALARMLOREG]{TIMG\_T\regindex{x}ALARMLO\_REG}、\hyperref[regdesc:TIMGTXALARMHIREG]{TIMG\_T\regindex{x}ALARMHI\_REG}、\hyperref[fielddesc:TIMGTXLOADLO]{TIMG\_T\regindex{x}\_LOAD\_LO} 和 \hyperref[fielddesc:TIMGTXLOADHI]{TIMG\_T\regindex{x}\_LOAD\_HI}。否则，上述寄存器应保持不变
        \item 置位 \hyperref[fielddesc:TIMGTXALARMEN]{TIMG\_T\regindex{x}\_ALARM\_EN} 重新使能报警
    \end{itemize}
    \item（最后一次报警时）关闭定时器
    \begin{itemize}
        \item 置位定时器在 \hyperref[fielddesc:TIMGTXINTCLR]{TIMG\_T\regindex{x}\_INT\_CLR} 的对应位清除中断
        \item 清零 \hyperref[fielddesc:TIMGTXEN]{TIMG\_T\regindex{x}\_EN} 关闭定时器
    \end{itemize}
\end{enumerate}

\subsection{通过 ETM 设置定时器用于周期性报警}
\label{sec:timg-conf-etm}

\begin{enumerate}
    \item 使能 ETM 模块的时钟
    \item 置位 \hyperref[fielddesc:TIMGTXAUTORELOAD]{TIMG\_T\regindex{x}\_AUTORELOAD} 使能自动重新加载；若由 ETM 任务重新加载计数器，则清零该字段
    \item 将 ETM 事件映射到相应的 ETM 任务（使用 ETM 事件触发对应的 ETM 任务）
    \begin{itemize}
        \item 如果 \hyperref[fielddesc:TIMGTXAUTORELOAD]{TIMG\_T\regindex{x}\_AUTORELOAD} 置 1，通过一个 ETM 通道映射 TG0\_EVT\_CNT\_CMP\_TIMER\regindex{x} 与 TG1\_EVT\_CNT\_CMP\_TIMER\regindex{x} 事件到 TG0\_TASK\_ALARM\_START\_TIMER\regindex{x} 与 TG1\_TASK\_ALARM\_START\_TIMER\regindex{x} 任务。
        \item 如果 \hyperref[fielddesc:TIMGTXAUTORELOAD]{TIMG\_T\regindex{x}\_AUTORELOAD} 置 0，除了映射 TG0\_EVT\_CNT\_CMP\_TIMER\regindex{x} 与 TG1\_EVT\_CNT\_CMP\_TIMER\regindex{x} 事件到 TG0\_TASK\_ALARM\_START\_TIMER\regindex{x} 与 TG1\_TASK\_ALARM\_START\_TIMER\regindex{x} 任务外，还需要用另一个 ETM 通道映射 TG0\_EVT\_CNT\_CMP\_TIMER\regindex{x} 与 TG1\_EVT\_CNT\_CMP\_TIMER\regindex{x} 事件到 TG0\_TASK\_CNT\_RELOAD\_TIMER\regindex{x} 与 TG1\_TASK\_CNT\_RELOAD\_TIMER\regindex{x} 任务。
    \end{itemize}
    \item 选择使能一个或两个 ETM 通道
    \item 将 \hyperref[fielddesc:TIMGETMEN]{TIMG\_ETM\_EN} 置 1 从而使能定时器组的 ETM 事件和任务
    \item 按照第 \ref{sec:timg-func-descr-simple-clock} 节 \textit{\nameref{sec:timg-func-descr-simple-clock}} 的第 1 步配置时基计数器
    \item 按照第 \ref{sec:timg-func-descr-timer-alarm} 节 \textit{\nameref{sec:timg-func-descr-timer-alarm}} 的第 2 步配置报警
    \item 通过 \hyperref[fielddesc:TIMGTXLOADLO]{TIMG\_T\regindex{x}\_LOAD\_LO} 和 \hyperref[fielddesc:TIMGTXLOADHI]{TIMG\_T\regindex{x}\_LOAD\_HI} 配置重新加载值
    \item 置位 \hyperref[fielddesc:TIMGTXEN]{TIMG\_T\regindex{x}\_EN} 开启定时器
    \item 置位 \hyperref[fielddesc:TIMGTXALARMEN]{TIMG\_T\regindex{x}\_ALARM\_EN} 开启报警
    \item 处理 TG0\_EVT\_CNT\_CMP\_TIMER\regindex{x} 与 TG1\_EVT\_CNT\_CMP\_TIMER\regindex{x}
    \begin{itemize}
        \item 当计数器达到报警值并触发警报时，硬件会产生 TG0\_EVT\_CNT\_CMP\_TIMER\regindex{x} 与 TG1\_EVT\_CNT\_CMP\_TIMER\regindex{x} 事件，同时报警被清除。
        \item 如果 \hyperref[fielddesc:TIMGTXAUTORELOAD]{TIMG\_T\regindex{x}\_AUTORELOAD} 为 1，硬件会自动将重新加载值写入计数器，TG0\_TASK\_ALARM\_START\_TIMER\regindex{x} 与 TG1\_TASK\_ALARM\_START\_TIMER\regindex{x} 重新使能报警。
        \item 如果 \hyperref[fielddesc:TIMGTXAUTORELOAD]{TIMG\_T\regindex{x}\_AUTORELOAD} 为 0，TG0\_TASK\_CNT\_RELOAD\_TIMER\regindex{x} 与 TG1\_TASK\_CNT\_RELOAD\_TIMER\regindex{x} 重新加载计数器，TG0\_TASK\_ALARM\_START\_TIMER\regindex{x} 与 TG1\_TASK\_ALARM\_START\_TIMER\regindex{x} 重新使能报警。
    \end{itemize}
    \item（最后一次报警时）关闭定时器
    \begin{itemize}
        \item 关闭用于映射定时器组事件和任务的 ETM 通道
        \item 将 \hyperref[fielddesc:TIMGETMEN]{TIMG\_ETM\_EN} 置 0
        \item 置位定时器在 \hyperref[fielddesc:TIMGTXINTCLR]{TIMG\_T\regindex{x}\_INT\_CLR} 的对应位清除中断
        \item 清零 \hyperref[fielddesc:TIMGTXEN]{TIMG\_T\regindex{x}\_EN} 关闭定时器
    \end{itemize}
\end{enumerate}

\subsection{慢速时钟频率计算}
\label{sec:rtc-calibration-procedure}

\begin{enumerate}
    \item 单次计算
    \begin{itemize}
        \item 设置 \hyperref[fielddesc:HPSYSCLKRSTTIMERGRP0TGRTCLKSRCSEL]{HP\_SYS\_CLKRST\_TIMERGRP0\_TGRT\_CLK\_SRC\_SEL} 选择需要计算频率的时钟，该字段设置值与时钟源具体对应关系参考章节 \ref{mod:resclk} \textit{\nameref{mod:resclk}}。

        \item 设置 \hyperref[fielddesc:HPSYSCLKRSTTIMERGRP0TGRTCLKDIVNUM]{HP\_SYS\_CLKRST\_TIMERGRP0\_TGRT\_CLK\_DIV\_NUM} 对所选时钟进行分频，以满足慢速时钟频率计算的前提条件。

        \item 设置 \hyperref[fielddesc:TIMGRTCCALIMAX]{TIMG\_RTC\_CALI\_MAX} 配置频率计算时间。

        \item 清零 \hyperref[fielddesc:TIMGRTCCALISTARTCYCLING]{TIMG\_RTC\_CALI\_START\_CYCLING} 选择单次校准模式，然后配置 \hyperref[fielddesc:TIMGRTCCALISTART]{TIMG\_RTC\_CALI\_START} 开启两个计数器。

        \item 等待 \hyperref[fielddesc:TIMGRTCCALIRDY]{TIMG\_RTC\_CALI\_RDY} 的值变为 1，读取 \hyperref[fielddesc:TIMGRTCCALIVALUE]{TIMG\_RTC\_CALI\_VALUE} 获取 XTAL\_CLK 计数器值，根据第 \ref{sec:rtc-calibration-desc} 节 \textit{\nameref{sec:rtc-calibration-desc}} 的公式计算慢速时钟的频率。
    \end{itemize}
    \item 周期性计算
    \begin{itemize}
        \item 设置 \hyperref[fielddesc:HPSYSCLKRSTTIMERGRP0TGRTCLKSRCSEL]{HP\_SYS\_CLKRST\_TIMERGRP0\_TGRT\_CLK\_SRC\_SEL} 选择需要计算频率的时钟。

        \item 设置 \hyperref[fielddesc:HPSYSCLKRSTTIMERGRP0TGRTCLKDIVNUM]{HP\_SYS\_CLKRST\_TIMERGRP0\_TGRT\_CLK\_DIV\_NUM} 对所选时钟进行分频，以满足慢速时钟频率计算的前提条件。

        \item 设置 \hyperref[fielddesc:TIMGRTCCALIMAX]{TIMG\_RTC\_CALI\_MAX} 配置频率计算时间。

        \item 使能 \hyperref[fielddesc:TIMGRTCCALISTARTCYCLING]{TIMG\_RTC\_CALI\_START\_CYCLING}，硬件将不间断进行频率计算过程。

        \item 只要 \hyperref[fielddesc:TIMGRTCCALICYCLINGDATAVLD]{TIMG\_RTC\_CALI\_CYCLING\_DATA\_VLD} 为 1，即表示 \hyperref[fielddesc:TIMGRTCCALIVALUE]{TIMG\_RTC\_CALI\_VALUE} 有效。
    \end{itemize}
    \item 超时\\
    如果 RTC 频率计算定时器的值超过 \hyperref[fielddesc:TIMGRTCCALITIMEOUTTHRES]{TIMG\_RTC\_CALI\_TIMEOUT\_THRES} 指定的阈值时，慢速时钟的计数器仍未完成计数，将置位 \hyperref[fielddesc:TIMGRTCCALITIMEOUT]{TIMG\_RTC\_CALI\_TIMEOUT} 标记计算超时。
\end{enumerate}




%%%%%%%%%%%%%%%%%%%%%%%%
%%% Register Summary %%%
%%%%%%%%%%%%%%%%%%%%%%%%

% >>> Update the sample text below and insert
%     the info on register summary into the subfile <<<

\newpage
\hypertarget{timg-reg-summ}{}
\section{寄存器列表}
\label{sec:timg-reg-summ}

本小节的所有地址均为相对于定时器组基地址的地址偏移量（相对地址），具体基地址请见章节 \ref{mod:sysmem} \textit{\nameref{mod:sysmem}} 中的表 \ref{tab:sysmem-base-address} \textit{\nameref{tab:sysmem-base-address}}。

请查看章节 \hyperref[glossary-access-types]{\textit{寄存器的访问类型}}，了解“访问”列缩写的含义。

\subfile{\modulefiles/17-TIMG-reg-summ__CN}


%%%%%%%%%%%%%%%%%
%%% Registers %%%
%%%%%%%%%%%%%%%%%

% >>> Update the sample text below and insert
%      the info on registers into the subfile <<<

\newpage
\section{寄存器}

本小节的所有地址均为相对于定时器组基地址的地址偏移量（相对地址），具体基地址请见章节 \ref{mod:sysmem} \textit{\nameref{mod:sysmem}} 中的表 \ref{tab:sysmem-base-address} \textit{\nameref{tab:sysmem-base-address}}。


\subfile{\modulefiles/17-TIMG-reg__CN}


\end{document}
